\documentclass[11pt]{article}
\usepackage[a4paper,margin=2.6cm]{geometry}
\usepackage{amsmath,amssymb,amsthm}
\newtheorem{lemma}{Lemma}
\newtheorem{proposition}[lemma]{Proposition}
\theoremstyle{remark}\newtheorem*{remark}{Remark}
\title{The quorum exponents:\\ \large what reduces, what the classical bounds miss, what remains}
\author{D.\ Joubert}
\date{3 September 2026}
\begin{document}
\maketitle
\begin{abstract}
\noindent
The every-scale quorum asks for two bounds on the bottom vector $v$ of
$Q_L$: $\delta_p=v^{\!\top}T_pv\le e^{-7p+c}$ and
$\kappa_p=\|T_pv-(v^{\!\top}T_pv)v\|\ge e^{-4p-c'}$, uniformly in the
window. Both quantities are pairings of $v$ against the prime tower
$T_p$. The leading term is the window autocorrelation of $v$ at lag
$\log p$. Bandlimited decay (Slepian, Bernstein, integration by parts)
gives powers of $1/\log p$, not $e^{-7p}$. The measured silence is the
decay of a near-kernel, not of a desert prolate. The hole shrinks to
one estimate: $\|T_pv\|$ itself is exponentially small in $p$. That
estimate is not proved. Nothing here is RH.
\end{abstract}

\section{Reduction}
$T_p$ is the contribution of the powers $p^k\le\mu$ to the prime-side
form,
\[
T_p(f,f)=\sum_{p^k\le\mu}(\log p)\,p^{-k/2}\,\Theta_f(k\log p),
\qquad
\Theta_f(y)=\int f(x)f(x+y)\,dx,
\]
with $\Theta_f$ supported in $[0,L]$, $L=\log\mu$, and $\Theta_f(L)=0$.
Then $\delta_p=T_p(v,v)$ and
$\kappa_p^2=\|T_pv\|^2-\delta_p^2$. The $k=1$ term dominates for
$p\ge5$:
\begin{equation}
\label{lead}
\delta_p=\frac{\log p}{\sqrt p}\,\Theta_v(\log p)+O\bigl(p^{-1}\Theta_v(2\log p)\bigr).
\end{equation}
So $\delta_p\le e^{-7p+c}$ is a pointwise bound on the autocorrelation
of one explicit vector at the points $\log p$.

\section{What the classical bounds give}
If $v$ were merely unit-norm in $PW_\tau$,
\[
|\Theta_v(y)|\le\min\bigl(1,\,C_N(\tau y)^{-N}\bigr)
\]
after $N$ integrations by parts on a smooth cutoff. At $y=\log p$ this
is a power of $1/\log p$. The Slepian that lives in the desert
$(-\gamma_1,\gamma_1)$ has
$\Theta(y)\sim\gamma_1\,\mathrm{sinc}(\gamma_1 y/\pi)$, which at
$p=7$ is $O(1/28)$, against a measured $\delta_7\sim10^{-16}$.
Duffin--Schaeffer on gaps loses the same exponential that
\emph{sampling-floor} already recorded for $c_L$.

None of these tools sees $e^{-7p}$. The exponent linear in $p$ is not
the Fourier decay of a function bandlimited to $\gamma_1$.

\section{What the measurements say instead}
At $\mu=11$, $-\ln\delta_p\approx7p-12$ for $p\ge5$, and later
$-\ln\delta_p\approx0.19\,s(\chi)\,p\log p$ across $\zeta,\chi_3,\chi_4$.
Both $\delta_p$ and $\kappa_p$ are exponentially small, so
$\|T_pv\|$ itself is exponentially small: $v$ is close to the kernel
of $T_p$, not just orthogonal to $T_pv$ in one direction. That is
stronger than silence. It means the near-kernel of $Q$ is already
almost orthogonal to every individual tower of a large prime.

This matches the leakage picture: the mass of $\hat v$ is expelled to
the band edge, and the pairing against a lag $\log p$ --- a frequency
modulation of width $1$ --- sees that edge through a rapidly
oscillating integral. Writing a bound that tracks the measured $7p$
requires a pointwise description of $\hat v$ on the edge, uniformly in
$L$. The leakage profile $\ln|\hat v(\gamma)|\approx-\tau(\omega_{\max}-\gamma)$
is measured, not proved, and $\tau$ itself moves with $\mu$.

\section{An attempt}
$\Theta_v$ is supported in $[0,L]$. The points $y=\log p$ lie
\emph{inside} the support. Exponential decay as $y\to\infty$ is the
wrong picture: there is no infinity. What is asked is a dip of size
$e^{-7p}$ at each arithmetic lag $\log p\in(0,L)$.

\paragraph{Near the edge.}
Bernstein plus $\Theta_v(L)=0$ gives
$|\Theta_v(\log p)|\le\tau\log(\mu/p)$. Useful only for $p=\mu^{1-o(1)}$.
At $p=7$, $\mu=11$, the bound is $O(1)$ against $10^{-16}$.

\paragraph{Contour shift.}
$F\in PW_\tau$ is entire of type $\tau$, so $\int F(\omega)\overline{F(\omega)}\,e^{i\omega y}\,d\omega$
can be shifted. The shift produces a factor $e^{-\delta y}$ and a
bound $e^{\tau\delta}$ on the new line --- net $e^{(-\delta+\tau)y}$ at
best, i.e.\ a power of $p$, and only after replacing $|F|^2$ on
$\mathbb R$ by $F(z)F(\bar z)$ with control we do not have on the
ground state. Super-exponential in $p=e^y$ never comes from a strip.

\paragraph{Fixed $p$, $L\to\infty$.}
The measured $\delta_p$ is stable from $\mu=8$ to $16$. The statement
uniform in $\mu$ is therefore: the Beurling extremal $v_L$ of the
zero set in $PW_{L/2}$ has $\Theta_{v_L}(\log p)$ bounded by $e^{-7p}$
for each fixed $p$, for all large $L$. That is a quantitative property
of one sampling-extremal function attached to $\Gamma$. It is the same
order of difficulty as a quantitative lower bound for $c_L$, already
recorded as out of reach of Duffin--Schaeffer / Bernstein
(\emph{sampling-floor}, \S4).

\section{Two lemmas that do hold}
\begin{lemma}[edge]
$\Theta_v(L)=0$, hence the tower of a prime sitting at the window edge
vanishes identically ($T_{11}\equiv0$ at $\mu=11$).
\end{lemma}
\begin{lemma}[crude]
$|\delta_p|\le(\log p)\,p^{-1/2}\|v\|_1^2$. In particular $\delta_p=O(p^{-1/2}\log p)$.
\end{lemma}
Lemma 2 is weaker than the $2\times2$ criterion by an exponential
factor, exactly as Bernstein was weaker than $c_L$ by $e^{L\gamma_1/2}$.

\section{Why $p\log p$, why a Gaussian}
The linear laws $-ln\delta\sim7p$ and $-ln\kappa\sim4p$ were local
slopes. In the weight $w=p\log p$ the measurements collapse
(\S87):
\[
-\ln\delta_p\approx0.19\,s(\chi)\,w,\qquad
-\ln\kappa_p\approx c(\mu)\,w^2\quad(p\ge5).
\]
One constant $0.19$ for $\zeta,\chi_3,\chi_4$ at four windows. $c(\mu)$
grows with the window ($0.005$--$0.03$). The $2\times2$ margin
$\ln(\kappa^2/\delta)\approx0.19 s\,w-2c\,w^2$ is then a downward
parabola: it peaks near $w\approx43$ and would vanish near $w\approx86$
($p\approx27$ at $\chi_3$, $\mu=30$), which is the measured finite
range of the single-vector certificate.

\paragraph{Why $w=p\log p$.}
$y=\log p$ is the additive lag at which $\Theta_v$ is evaluated.
$e^y=p$ is the multiplicative place of the same prime. The product
$y\,e^y=p\log p$ is the phase-space area of that prime in the
window --- lag times place. Neither factor alone collapsed the three
$L$-functions: $p$ gave the old $7p$ (local), $\log p$ is too slow,
$(\log p)/\sqrt p$ is the bare tower weight and goes the wrong way.
Stirling's $p\log p-p$ is the same area: $\Theta_v(\log p)\sim1/p!$
would produce exactly this leading term. That reading is a mnemonic,
not a derivation. What is structural is that the form multiplies an
additive pairing by an Euler place, and the invariant of that product
is $w$.

\paragraph{Why a Gaussian in $w$.}
Silence is linear in $w$: one cumulant, intensive in the speed $s$.
Coupling is the remainder of $T_pv$ orthogonal to $v$. A quadratic
cost in the same action is the next cumulant. Then $\kappa\sim
e^{-c(\mu)w^2}$ is Gaussian in $w$, and $c$ may depend on the window
because the orthogonal of $v$ grows with $N$. This is maximum-entropy
bookkeeping, not a spectral theorem. It does explain, without a new
parameter, why the margin is a parabola and why a single vector
stops certifying around $p=29$.

A check that would kill the reading: a fourth $L$-function whose
$-ln\delta_p/(s\,w)$ is not $0.19$, or a window at which $-ln\kappa$
is linear in $w$ rather than quadratic. $\chi_5$ or $\mu=38$ on
$\zeta$ would do.

\section*{Status}
Tried. The three routes --- decay at infinity, contour shift, edge
continuity --- miss $e^{-7p}$ for the same structural reason: the lag
is interior and the smallness is a dip of the Beurling extremal at
arithmetic points. The measured laws live in the phase-space weight
$w=p\log p$; the Gaussian in $w$ for $\kappa$ is the second cumulant
of that action. Neither law is proved. The hole is the same hole as
the quantitative floor. No RH.
\end{document}
